\section{The hairpin geometry is self-contradictory}
\label{sec:hairpin}

\subsection{What the cascade requires}

The charging water enters the borehole at $\Tm + \Delta T_{4C\text{-}M}$ and
glides $\Delta T_{3C\text{-}2C}$ as it gives up heat, so with the default
values $\Delta T_{4C\text{-}M} = \SI{10}{\kelvin}$ and
$\Delta T_{3C\text{-}2C} = \SI{55}{\kelvin}$ it enters at \SI{160}{\degC} and
leaves at \SI{105}{\degC}. The cascade exists so that the melting temperature
tracks the water: each layer sits a fixed approach below the local fluid
temperature, and the span of melting temperatures matches the glide.

In a hairpin the water traverses two legs in series. Over the second of them
it is already below the top melting temperature --- indeed, with the figures
above, the return leg runs
\begin{equation}
\Tm - (T_{4C} - \Delta T_{3C\text{-}2C}) = 150 - 105 = \SI{45}{\kelvin}
\end{equation}
\emph{below} the melting point. If the cascade were graded by depth, so that
both legs at a given depth saw the same $\Tm$, the return leg would freeze what
the down leg had just melted. The model therefore graded the cascade along the
developed path $z \in [0, 2\Lw]$.

\subsection{What path grading costs}

Path grading is self-consistent, and it is what THUMS did. But it requires two
distinct melting temperatures to coexist at the same depth. With a linear
grade over the path, the difference between the down leg at path position $d$
and the return leg at $2\Lw - d$ is
\begin{equation}
\Delta\Tm(d) = \frac{\text{span}}{2\Lw}\,\big[(2\Lw - d) - d\big]
  = \text{span}\left(1 - \frac{d}{\Lw}\right),
\end{equation}
which is the full span of \SI{48.9}{\kelvin} at the wellhead, falling linearly
to zero at the bottom, with a depth-average of \SI{24.4}{\kelvin}.

These two bodies of material are separated by centimetres of \emph{contiguous}
phase-change material, with no wall between them. Treating the legs as two
parallel cylinders of radius $r_e$ with centre spacing $d_{cc}$ in a medium of
conductivity $k_p$, the conduction shape factor per unit length is
\begin{equation}
\frac{S}{L} = \frac{2\pi}{\operatorname{arcosh}
  \!\left(\dfrac{d_{cc}^{2} - 2r_e^{2}}{2r_e^{2}}\right)},
\qquad
\qp_{\text{leg-leg}} = k_p \frac{S}{L}\, \overline{\Delta\Tm}.
\end{equation}

\begin{table}[h]
\centering
\begin{tabular}{lrrrr}
\toprule
hairpins & legs & $d_{cc}$ & $\qp_{\text{leg-leg}}$ & field total, \% of duty \\
\midrule
1 & 2 & \SI{126}{\milli\metre} & \SI{28}{\watt\per\metre} & \SI{43}{\kilo\watt} \ (\SI{36}{\percent}) \\
2 & 4 & \SI{89}{\milli\metre}  & \SI{41}{\watt\per\metre} & \SI{62}{\kilo\watt} \ (\SI{52}{\percent}) \\
3 & 6 & \SI{73}{\milli\metre}  & \SI{60}{\watt\per\metre} & \SI{92}{\kilo\watt} \ (\SI{77}{\percent}) \\
\bottomrule
\end{tabular}
\caption{Parasitic conduction between the legs of a hairpin, evaluated at the
depth-average temperature difference. This term is absent from the THUMS
model. Note its direction: it \emph{grows} with the number of hairpins, which
the v0.16 parametric study ranked as the single most beneficial design
variable, with an effect of $+\SI{33.7}{\percent}$ on
$\varepsilon_{\mathrm{cyc}}\eta_{\mathrm{RTE}}$.}
\label{tab:shortcircuit}
\end{table}

The heat does not leave the store --- it moves from hot material to cold
material within it --- so this is an exergy destruction and a degradation of
the cascade rather than an energy loss. That does not make it benign. The
cascade is the mechanism by which the store matches a gliding source, and a
parasitic of \SI{77}{\percent} of duty flowing directly down the temperature
grade is a short circuit across the model's central idea.

\subsection{Two further consequences}

\paragraph{Path-graded material is not pourable.} Two melting temperatures at
one depth, in separate radial cells, requires radially segregated cartridges
rather than material placed in layers from the bottom up. This is a
manufacturing requirement that the geometry implies and that was never stated.

\paragraph{No depth-dependent boundary condition is possible.} With
$L_t = 2\Lw$ and $z$ a path length, the map $z \mapsto \text{depth}$ is
two-to-one: each depth is occupied by the down leg at $z = d$ and the return
leg at $z = 2\Lw - d$. A formation temperature that varies with depth would
have two segments claiming it, with no principled way to divide the loss
between them. \emph{This is why the two changes in BOREAS are one change}: the
conducting wall cannot be attached to the hairpin geometry at all.

